\subsection{Сложение: сначала приводим к одному масштабу}
Складывать значащие части напрямую можно только при одинаковых порядках.
Как $3\text{ м}+20\text{ см}$ требуют общих единиц, так и
$M_a2^{E_a}+M_b2^{E_b}$ требуют общего масштаба.

\paragraph{Алгоритм для конечных операндов.}
\begin{enumerate}
    \item Восстановить знак, значащую часть и истинный порядок.
    Для субнормальных использовать ведущий 0 и масштаб $2^{-126}$
    в binary32, а не $2^{-127}$.
    \item Выбрать больший порядок. Значащую часть меньшего операнда
    сдвинуть вправо на разность порядков, сохраняя информацию о хвосте.
    \item При одинаковых знаках сложить модули; при разных --- вычесть
    меньший модуль из большего и взять знак большего по модулю операнда.
    \item Нормализовать результат с коррекцией порядка. Отдельно
    обработать ноль и достижение нижней границы нормальных чисел.
    \item Округлить до точности формата. Если округление породило перенос,
    ещё раз скорректировать порядок. Проверить границы диапазона и упаковать поля.
\end{enumerate}

\subsubsection*{Точное сложение $6.5+0.75$ в binary32}
Представим операнды и приведём их к порядку $E=2$:
\[
\begin{aligned}
6.5 &= 1.101_2\cdot2^2,\\
0.75 &= 1.1_2\cdot2^{-1}=0.0011_2\cdot2^2.
\end{aligned}
\]
Вторую значащую часть сдвинули вправо на $2-(-1)=3$ позиции.
Само число не изменилось, потому что одновременно изменён масштаб.
\[
\begin{array}{r}
1.1010_2\\
+\ 0.0011_2\\
\hline
1.1101_2
\end{array}
\qquad
1.1101_2\cdot2^2=7.25.
\]
Результат уже нормализован. Все ненулевые биты помещаются в 24 значащих
разряда; округление ничего не меняет. Слово binary32: \texttt{0x40E80000}.

\paragraph{Если сумма вышла за $[1,2)$.}
При $1.5+1.5$ получаем $1.1_2+1.1_2=11.0_2$.
Перенос не отбрасывают: $11.0_2\cdot2^0=1.10_2\cdot2^1=3$.
Сдвиг вправо на одну позицию сопровождается увеличением порядка на 1.

\subsection{Округление и биты guard, round, sticky}
Простое отбрасывание хвоста --- \textbf{усечение}, а не округление RN.
После приведения результата к нужному масштабу выделим:
\begin{itemize}
    \item $L$ --- младший \emph{сохраняемый} бит значащей части;
    \item $G$ (guard) --- первый отбрасываемый бит;
    \item $R$ (round) --- второй отбрасываемый бит;
    \item $S$ (sticky) --- логическое ИЛИ всех оставшихся битов хвоста.
\end{itemize}
Sticky хранит ответ «есть ли дальше хотя бы одна единица», а не значение
одного следующего двоичного разряда. Для RN модуль сохраняемой части
увеличивают на единицу её младшего разряда тогда и только тогда, когда
\[
G\mathbin{\land}(R\mathbin{\lor}S\mathbin{\lor}L)=1.
\]
Если $G=0$, хвост меньше половины шага. При $G=1$ и $R=S=0$
получаем ровно половину: решение зависит от чётности $L$.
При $G=1$ и ненулевом $R$ или $S$ хвост больше половины.
Для отрицательного результата правило применяется к модулю, затем
возвращается знак. GRS оцениваются относительно окончательной границы
сохраняемых разрядов; это не рецепт сдвигать sticky как обычный бит
при любой последующей нормализации.

\begin{center}
\begin{tabular}{|l|c|c|l|}
\hline
Точное значение, $p=4$ & $L$ & $GRS$ & Результат \\
\hline
$1.01001_2$ & 0 & 010 & $1.010_2$ \\
$1.0101_2$ & 0 & 100 & $1.010_2$ \\
$1.010101_2$ & 0 & 101 & $1.011_2$ \\
$1.0111_2$ & 1 & 100 & $1.100_2$ \\
$1.1111_2$ & 1 & 100 & $1.000_2\cdot2^1$ \\
\hline
\end{tabular}
\end{center}
В последней строке округление вызвало перенос: $1.111_2+0.001_2=10.000_2$.
Нужно заново нормализовать число, а не потерять старшую единицу.

\paragraph{Сложение binary32 чуть выше середины.}
Возьмём точно представимые $a=1$ и
$b=1.01_2\cdot2^{-24}=2^{-24}+2^{-26}$. После выравнивания
точная сумма равна $1+2^{-24}+2^{-26}$.
Храним 23 дробных бита: они все нулевые, $L=0$; далее
$G=1$, $R=0$, $S=1$. Поэтому
\[
\operatorname{RN}(a+b)=1+2^{-23}
\quad (\texttt{0x3F800001}).
\]
Если забыть sticky, хвост ошибочно станет «ровно половиной» и получится 1.
Если бы слагаемое было ровно $2^{-24}$, результат действительно был бы 1.

\subsection{Вычитание: нормализация влево}
Для $1.125-1$ оба порядка равны нулю:
\[
1.001_2-1.000_2=0.001_2=1.000_2\cdot2^{-3}=0.125.
\]
Ведущая единица оказалась на три позиции правее. Сдвигаем значащую
часть влево на три позиции и уменьшаем порядок на 3. Слово результата
binary32: \texttt{0x3E000000}. \textbf{Вычитание здесь точное}:
исходные значения и разность представимы без округления.
Нули, добавленные при нормализации, не восстанавливают ранее утраченную
информацию; в этом примере её просто не теряли.

Если модули равны и знаки противоположны, получается ноль, который
нельзя нормализовать до ведущей единицы. Если при сдвиге достигается
минимальный нормальный порядок, дальнейший результат представляют
как субнормальный с ведущим 0. Разность с отрицательным знаком
обрабатывается так же по модулю.

\subsection{Отмена близких чисел и погрешность исходных данных}
Пусть вместо точных $x,y$ известны $\hat x=x+\Delta x$ и
$\hat y=y+\Delta y$. Даже \emph{точное} вычитание приближений даёт
\[
 (\hat x-\hat y)-(x-y)=\Delta x-\Delta y.
\]
Абсолютная ошибка разности не превосходит
$|\Delta x|+|\Delta y|$, но относительно маленького $|x-y|$ может
оказаться большой. Общие ведущие разряды сокращаются (cancellation),
а ошибки исходных данных остаются. Это объясняет возможную
катастрофическую потерю относительной точности.

\paragraph{Проверенный пример из десятичных данных.}
Пусть $x=1.23456789$, $y=1.23456780$. Точная разность равна $9\cdot10^{-8}$.
В binary32 с RN получаем
\[
\begin{aligned}
\hat x &= 1.2345678806304931640625,\\
\hat y &= 1.23456776142120361328125.
\end{aligned}
\]
Слова: \texttt{0x3F9E0652} и \texttt{0x3F9E0651}. Их разность
\[
\hat x-\hat y=2^{-23}\approx1.1920929\cdot10^{-7}
\]
представляется \emph{точно}. Однако относительно исходной разности
ошибка составляет около $32.45\%$. Она возникла при представлении
операндов, а вычитание сделало её заметной. Поэтому утверждение
«вычитание близких чисел всегда неточно» неверно.

\paragraph{Что может помочь.}
Использовать более высокую точность \emph{с самого начала}, оценивать
качество входных данных, выбирать менее чувствительную формулу.
Например, при малом ненулевом $x>-1$ выражение
\[
\sqrt{1+x}-1=\frac{x}{\sqrt{1+x}+1}
\]
лучше вычислять в правой форме: она избегает вычитания двух чисел,
близких к 1. Это пример выбора устойчивой формулы, а не универсальное
решение всех численных проблем. Преобразование уже округлённого
\texttt{float} в \texttt{double} не восстановит потерянные цифры.

\subsection{Умножение: значащие части умножаем, порядки складываем}
Для ненулевых нормальных операндов
\[
xy=(-1)^{s_x\oplus s_y}(M_xM_y)2^{E_x+E_y}.
\]
Знак определяется XOR знаков. Так как $1\le M_x,M_y<2$,
то $1\le M_xM_y<4$. Если произведение не меньше 2, сдвигаем его
вправо и увеличиваем порядок на 1. Затем округляем, учитывая
возможный новый перенос. Складываются \emph{истинные} порядки:
для хранимых полей перед нормализацией получается
$e_x+e_y-\mathrm{bias}$, а не $e_x+e_y$.

\paragraph{Пример $1.5\times2.5$.}
\[
\begin{aligned}
1.5 &= 1.1_2\cdot2^0,& 2.5 &= 1.01_2\cdot2^1,\\
1.1_2\cdot1.01_2 &=1.1_2+0.011_2=1.111_2.
\end{aligned}
\]
Следовательно, результат $1.111_2\cdot2^1=3.75$ точен в binary32.
Для сравнения, $1.5\times1.5=10.01_2\cdot2^0=1.001_2\cdot2^1=2.25$:
здесь требуется нормализация. В общем случае произведение двух
$p$-разрядных значащих частей может потребовать до $2p$ разрядов
до окончательного округления.

\subsection{Деление: значащие части делим, порядки вычитаем}
При конечном ненулевом делителе и нормальных операндах
\[
\frac{x}{y}=(-1)^{s_x\oplus s_y}\frac{M_x}{M_y}\,2^{E_x-E_y}.
\]
Частное значащих частей лежит в $(0.5,2)$. Если оно меньше 1,
сдвигаем влево и уменьшаем порядок. Деление может породить
бесконечную двоичную дробь; остаток деления помогает определить,
есть ли ненулевой хвост для округления.

\paragraph{Точное деление $3/1.5$.}
\[
\frac{1.1_2\cdot2^1}{1.1_2\cdot2^0}=1.0_2\cdot2^1=2.
\]
\paragraph{Неточное деление $1/3$.}
\[
\frac{1.0_2\cdot2^0}{1.1_2\cdot2^1}
=0.101010\ldots_2\cdot2^{-1}
=1.010101\ldots_2\cdot2^{-2}.
\]
После 23 дробных битов в binary32 $L=0$, $G=1$, $R=0$, $S=1$.
Округляем вверх по модулю. Получаем \texttt{0x3EAAAAAB}, то есть
\[
\hat x=\frac{11184811}{33554432}
\approx0.3333333432674408.
\]
Абсолютная ошибка равна $1/100663296\approx9.93\cdot10^{-9}$,
относительная --- $2^{-25}\approx2.98\cdot10^{-8}$.

\subsection{Поглощение и неассоциативность}
\paragraph{Поглощение.}
В binary32 справа от $2^{24}=16777216$ шаг равен 2.
Точная сумма $2^{24}+1$ находится посередине двух соседей.
RN выбирает $2^{24}$ с чётным младшим разрядом.
Малое слагаемое участвовало в операции, но \emph{округлённый результат}
не изменился. Нельзя заранее отбросить любой операнд, меньший одного
шага: значение больше половины шага уже может изменить результат.

\paragraph{Скобки имеют значение.}
Для $a=2^{24}$, $b=1$, $c=-2^{24}$, с округлением каждого действия
в binary32:
\[
\begin{aligned}
\operatorname{fl}(\operatorname{fl}(a+b)+c)&=0,\\
\operatorname{fl}(a+\operatorname{fl}(b+c))&=1.
\end{aligned}
\]
Во второй строке $b+c=-16777215$ представимо точно: на стороне
меньших модулей шаг ещё равен 1.
\begin{lstlisting}[language=C,basicstyle=\ttfamily\small]
float a = 16777216.0f, b = 1.0f, c = -16777216.0f;
float ab = a + b, bc = b + c;
float left = ab + c, right = a + bc;
printf("%.0f %.0f\n", (double)left, (double)right);
// 0 1
\end{lstlisting}
В реальных суммах помогают накопление в \texttt{double}, попарное
суммирование или компенсация ошибки. Для неотрицательных слагаемых
часто полезно начинать с малых, но универсальной гарантии такой порядок
не даёт. Параллельное суммирование может менять порядок операций.
Опция \texttt{-ffast-math} разрешает преобразования, несовместимые
с использованной здесь пошаговой моделью.

\subsection{Границы диапазона и особые значения}
Максимальное конечное binary32 равно
$(2-2^{-23})2^{127}$ (\texttt{0x7F7FFFFF}). При RN его умножение на 2
даёт $+\infty$. Рост порядка не означает появления дополнительных
значащих битов: диапазон и точность --- разные характеристики.

\paragraph{Постепенный переход к нулю.}
Минимальное положительное normal равно $2^{-126}$, subnormal --- $2^{-149}$.
В субнормальном диапазоне шаг постоянен, $2^{-149}$, поэтому при
приближении к нулю относительная точность ухудшается.
\[
\begin{aligned}
2^{-126}/2 &= 2^{-127} &&\text{точный субнормальный результат},\\
\operatorname{RN}(2^{-149}/2)&=+0 &&\text{половина шага: к чётному нулю},\\
\operatorname{RN}(3\cdot2^{-150})&=2^{-148} &&\text{другая середина: к двум шагам}.
\end{aligned}
\]
Получение субнормального значения само по себе не означает ошибку
округления. При стандартной обработке флаг underflow связан с малостью
\emph{и неточностью}; точный результат $2^{-127}$ его не требует.
Аппаратные режимы FTZ/DAZ, принудительно заменяющие маленькие числа
нулём, в примерах не используются.

При обычной обработке IEEE 754 без прерывания на исключениях:
\begin{center}
\begin{tabular}{|l|l|}
\hline
Операция или проверка & Результат \\
\hline
$1/(+0)$, $1/(-0)$ & $+\infty$, $-\infty$ \\
$0/0$, $\infty-\infty$, $0\cdot\infty$ & NaN \\
$+0=-0$ & Истина при числовом сравнении \\
NaN $=$ NaN & Ложь \\
NaN $\ne$ NaN & Истина \\
\hline
\end{tabular}
\end{center}
Для проверки в C используйте \texttt{isnan}, \texttt{isfinite},
\texttt{signbit} из \texttt{<math.h>}. Упорядоченные сравнения с NaN
ложны. Флаги исключений плавающей арифметики --- не целочисленные
\texttt{CF}, \texttt{OF}, \texttt{ZF}, \texttt{SF} из лекции №2.

\subsection{Абсолютная, относительная ошибка и допуски}
Для точного $x$ и приближения $\hat x$:
\[
\Delta_{\rm abs}=|\hat x-x|,\qquad
\Delta_{\rm rel}=\frac{|\hat x-x|}{|x|}\quad(x\ne0).
\]
Абсолютная ошибка имеет единицы исходной величины; относительная ---
безразмерна. Вблизи нуля относительная ошибка часто неудобна.
Оценка ошибки по известному точному значению и проверка близости двух
приближений --- разные задачи: близость не доказывает правильности.

\paragraph{Правило близости для конечных значений.}
Выбрав неотрицательные $\delta_{\rm abs}$ и $\delta_{\rm rel}$, проверяют
\[
|a-b|\le\max\bigl(\delta_{\rm abs},
\delta_{\rm rel}\max(|a|,|b|)\bigr).
\]
Абсолютный допуск работает у нуля, относительный масштабируется
вместе с числами. Допуски выбирают по точности входных данных,
погрешности метода и требованиям к результату.

\begin{itemize}
    \item Температуры $20$ и $20.004$ при допустимом отклонении
    $0.01$ градуса считаем близкими: $\delta_{\rm abs}=0.01$,
    $\delta_{\rm rel}=0$. Это условие задачи, не свойство \texttt{double}.
    \item $10^8$ и $10^8+5$ при $\delta_{\rm rel}=10^{-6}$ близки:
    разрешённая разность около 100.
    \item $0$ и $10^{-10}$ не проходят только относительную проверку
    с $\delta_{\rm rel}=10^{-6}$. При $\delta_{\rm abs}=10^{-9}$ проходят.
\end{itemize}

\paragraph{Учебная функция на C.}
Здесь допуски должны быть конечными, $\delta_{\rm abs}\ge0$,
$0\le\delta_{\rm rel}<1$. Нечисловые и бесконечные аргументы
считаются недопустимыми для этой проверки близости.
\begin{lstlisting}[language=C,basicstyle=\ttfamily\small]
#include <stdbool.h>
#include <math.h>

bool close_finite(double a, double b,
                  double atol, double rtol) {
    if (!isfinite(a) || !isfinite(b) ||
        !isfinite(atol) || !isfinite(rtol) ||
        atol < 0 || rtol < 0 || rtol >= 1)
        return false;
    double scale = fmax(fabs(a), fabs(b));
    return fabs(a - b) <= fmax(atol, rtol * scale);
}
\end{lstlisting}
Функция предполагает обычный IEEE-режим без ловушек: если разность
конечных аргументов переполнится до бесконечности, условие даст ложь,
что корректно при указанных конечных допусках и $\delta_{\rm rel}<1$.
При сравнении очень близко к границе допуска сами арифметические
действия проверки тоже округляются.

Равенство \texttt{==} уместно, когда требуется точное совпадение
представимых значений, например проверка на ноль, который алгоритм
явно присваивает. Приближённая близость не транзитивна: при абсолютном
допуске 1 числа 0 и 0.75 близки, 0.75 и 1.5 близки, но 0 и 1.5 --- нет.
Поэтому такую проверку нельзя механически использовать вместо отношения
равенства для ключей таблиц или порядка сортировки.

\subsection*{Итоги и вопросы для самопроверки}
\begin{enumerate}
    \item Зачем перед сложением приводить числа к общему порядку?
    \item Что отличает $GRS=100$ от $GRS=101$ при $L=0$?
    \item Почему $1.125-1$ вычисляется точно, хотя требует сдвига влево?
    \item На каком этапе потеряна информация в примере с $1.23456789$?
    \item Почему $2^{24}+1$ не меняет binary32-результат?
    \item Почему машинный эпсилон не является универсальным допуском?
\end{enumerate}
\textbf{Краткие ответы:} общий вес разрядов; середина против значения
выше середины; операнды и разность представимы; при округлении входов;
середина между соседями с выбором чётного; допуск зависит от масштаба,
данных и требований задачи.

\newpage
\subsection{Дополнительно: полные раскладки и расширенные случаи}
Этот раздел не входит в основной хронометраж. Он дополняет разборы
и позволяет восстановить вычисления после лекции.

\paragraph{Полное сложение $6.5+0.75$.}
Поля входов:
\begin{center}
\texttt{0 10000001 10100000000000000000000}\\
\texttt{0 01111110 10000000000000000000000}.
\end{center}
Восстанавливаем ведущие единицы, выравниваем к $E=2$:
\begin{verbatim}
  1.10100000000000000000000
+ 0.00110000000000000000000
---------------------------
  1.11010000000000000000000
\end{verbatim}
$GRS=000$; $e=2+127=129$. Ответ:
\begin{center}
\texttt{0 10000001 11010000000000000000000}.
\end{center}

\paragraph{Полное вычитание $1.125-1$.}
Входы \texttt{0x3F900000}, \texttt{0x3F800000} имеют $e=127$.
\begin{verbatim}
  1.00100000000000000000000
- 1.00000000000000000000000
---------------------------
  0.00100000000000000000000
\end{verbatim}
После сдвига влево на 3: $M=1$, $E=-3$, $e=124$.
Ответ:
\begin{center}
\texttt{0 01111100 00000000000000000000000}.
\end{center}

\paragraph{Округление точной суммы и округление входов --- разные задачи.}
Сохраним старый учебный пример: точная сумма $1+1.001_2\cdot2^{-3}$
равна $1.001001_2=1.140625$. Её округление до $p=3$ даёт $1.01_2=1.25$,
поскольку середина между 1 и 1.25 равна 1.125.
Но второй операнд сам \emph{не представим} при $p=3$.
Если сначала округлить его до $1.00_2\cdot2^{-3}=0.125$,
сумма окажется ровно на середине и RN даст 1.
Нельзя выдавать округление точной суммы произвольных входов за
операцию над заранее округлёнными операндами учебного формата.

\paragraph{Числа на границе normal/subnormal.}
Максимальное субнормальное binary32 равно $2^{-126}-2^{-149}$,
слово \texttt{0x007FFFFF}. Следующее число --- минимальное normal,
\texttt{0x00800000}. Их разность $2^{-149}$ представима точно,
слово \texttt{0x00000001}: разрыва в сетке у этой границы нет.

\paragraph{Эпсилон и единица округления.}
Для RN часто используют $u=2^{-p}=\varepsilon/2$ (unit roundoff).
Если точный ненулевой результат находится в нормальном диапазоне
и нет переполнения, стандартная оценка имеет вид
\[
\operatorname{fl}(z)=z(1+\delta),\qquad |\delta|\le u.
\]
У субнормальных такая равномерная относительная граница не действует;
абсолютная ошибка округления к ближайшему не превосходит
половины субнормального шага. В разных книгах термин «машинный эпсилон»
может означать $\varepsilon$ или $u$, поэтому определение важно.

\paragraph{Другие режимы и FMA.}
Помимо RN существуют округления к нулю, к $+\infty$ и к $-\infty$.
При переполнении они могут давать максимальное конечное значение,
а не бесконечность. Все основные примеры лекции используют RN.
Операция FMA вычисляет $ab+c$ с \emph{одним} окончательным округлением,
а отдельные умножение и сложение --- с двумя. Например, в binary32
$a=1+2^{-23}$, $b=1-2^{-23}$, $c=-1$:
\[
\operatorname{RN}(\operatorname{RN}(ab)+c)=0,
\qquad \operatorname{RN}(ab+c)=-2^{-46}.
\]
Поэтому разрешённое компилятору объединение операций также может менять
результат. Стандартная функция C \texttt{fmaf} задаёт именно fused-операцию.

\paragraph{Источники.}
IEEE 754-2019~\cite{ieee-754}; статья Д.~Голдберга о
погрешностях~\cite{goldberg-floating-point}; проект стандарта C
N1570~\cite{c11-draft}, в частности приложение F.
